\documentclass[twoside,12pt]{dscproblemset}

\usepackage[margin=1in]{geometry}
\usepackage{amsmath}
\usepackage{amssymb,amsthm}
\usepackage{fancyhdr}
\usepackage{comment}
\usepackage{nicefrac}
\usepackage{needspace}
\usepackage{tikz}
\usepackage{forest}
\usepackage{pgfplots}
\usepackage{booktabs}
\usepackage[outputdir=build]{minted}
\usetikzlibrary{quotes,angles,positioning,arrows.meta}
\usetikzlibrary{calc}

% configuration
% ------------------------------------------------------------------------------

% page headers only on odd pages
\pagestyle{fancy}
\fancyhead{}
\fancyhead[RO]{PID or Name: \rule{3in}{.5pt}}
\renewcommand{\headrulewidth}{0pt}

% commands
% --------------------------------------------------------------------------------------

\newcommand{\python}[1]{\mintinline{python}{#1}}

\newcommand{\newpageinexam}{
    \ifshowsoln
    \else
        \newpage
    \fi
}

% --------------------------------------------------------------------------------------


\begin{document}

\examfrontpage{
    % course
    DSC 40B
}{
    % exam name
    Redemption Midterm 02
}{
    % exam date
    December 15, 2023
}{
    % instructions
    \textbf{Instructions:}
    \begin{itemize}
        \item Your exam version should be different from those to your left/right.
        \item Each problem is worth one point unless stated otherwise.
        \item Write your solutions to the following problems in the spaces provided.
        \item No calculators are permitted, but some pages of notes are.
        \item Write your name or PID at the top of each sheet in the space provided.
        \item Unless otherwise stated, you should assume that a node's neighbors are produced in ascending order
            of their label.
        \item If it is not stated, you may not assume that a graph is directed or undirected, connected or
            disconnected, etc.
    \end{itemize}
    \vspace{1em}
}

\begin{probset}

% - Hashing: 1
% - Graph theory: 5
%     - Number of nodes/edges: 2
%     - Connected components: 1
%     - Paths, reachability: 2
% - Graph representations: 2
% - BFS: 5
%     - Run a BFS (shortest paths): 1
%     - Property of the queue: 2
%     - Amortized time of modification, or time complexity of BFS on a special graph: 1-2
%     - Property of shortest paths: (1?)
% - DFS: 5
%     - Run a DFS (to find start and finish times): 1
%     - Property of start and finish times: 2
%     - Amortized time of modification, or time complexity of DFS on a special graph: 1
%     - Topological sorting: 1
% - Bellman-Ford: 1
%     - Run it: 1
%     - Loop invariant: 1
% - Dijkstra: 1
%     - Run it:

        \inputproblem{01}
        \inputproblem{02}
        \inputproblem{03}
        \inputproblem{04}
        \newpageinexam
        \inputproblem{05}
        \inputproblem{06}
        \inputproblem{07}
        \inputproblem{08}
        \inputproblem{09}
        \inputproblem{10}
        \inputproblem{11}
        \inputproblem{12}
        \newpageinexam
        \inputproblem{13}
        \inputproblem{14}
        \inputproblem{15}
        \inputproblem{16}
        \inputproblem{17}
        \inputproblem{18}
        \inputproblem{19}
        \inputproblem{20}

\end{probset}


\vfill
\begin{center}
    \textbf{Before turning in your exam, please check that your name is on every page.}
\end{center}


\ifshowsoln
\else
    % \scratchpage{}
\fi

\end{document}
